\section{Ansible}
Configures provisioned machines (install/configure software, deploy code); \textit{after} provisioning. Idempotent (no-op if already in desired state), agentless, push-based, no statefile/locking. Needs Python on host \& targets. Other CM tools: \textbf{Puppet}, \textbf{Chef}, \textbf{CFEngine}, \textbf{Juju}. \\
\textbf{TF vs Ansible:} TF = provisioning, declarative, state, idempotent. Ansible = CM, \textit{hybrid} (imperative tasks, declarative modules), no lifecycle, idempotency \textit{optional}. Both: templating, agentless, OSS. Often combined.\\
\textbf{Execution:} over SSH, copies generated Python module to \texttt{\textasciitilde/.ansible/tmp} on target, runs with remote Python, returns JSON on stdout, deletes. Nothing stays installed. \textbf{Exceptions:} \texttt{raw} module runs cmd via SSH (no Python needed, e.g. bootstrap); \texttt{connection: local} runs on controller (network devices); Windows via \textbf{WinRM}.
\subsection{Initial Setup \& Inventory}
\texttt{ssh-keygen -t ed25519 -C "x@y"} on controller, \texttt{ssh-copy-id user@server}, test with \texttt{ansible all -m ping}. \\
\textbf{Inventory} (INI or YAML) lists managed hosts grouped, referenced by play's \texttt{hosts:}. \texttt{ansible-inventory -i inv --list} verifies it.
\begin{lstlisting}
[web]
web1.example.com
web2.example.com ansible_host=10.0.0.5
[web:vars]
ansible_user=deploy
\end{lstlisting}
\subsection{Playbooks, Plays \& Tasks}
YAML file = playbook (1..n plays); play = hosts + tasks; task = one module call. Run: \texttt{ansible-playbook -i inv playbook.yml}. Flags: \texttt{--check} (dry run), \texttt{--list-tasks} (preview), \texttt{--tags}/\texttt{--skip-tags}, \texttt{--ask-vault-pass}.
\begin{lstlisting}[language=yaml]
- name: Example Playbook
  hosts: web
  become: true             # sudo
  gather_facts: false      # skip setup (faster)
  force_handlers: true     # run handlers even if a task fails
  vars:
    packages: [nginx, curl]
    secret: "{{ vault_password }}"
  roles:
    - myrole
  tasks:
    - name: Install packages    # list-as-option > loop (faster, deps OK)
      ansible.builtin.apt:      # FQCN: namespace.collection.module
        name: "{{ packages }}"  # alternative for apt is yum
        state: present
      notify: restart nginx
    - name: Render config
      ansible.builtin.template: { src: app.conf.j2, dest: /etc/app.conf }
      when: enable_service
      tags: config
  handlers:
    - name: restart nginx
      ansible.builtin.service: { name: nginx, state: restarted }
\end{lstlisting}
\textbf{Loop anti-pattern:} a per-item \texttt{loop:} runs the module once per package -- slower \& can break on interdependencies; prefer the whole list in one \texttt{name:} (single transaction, as above):
\begin{lstlisting}[language=yaml]
- name: suboptimal              # one apt/yum call per item
  ansible.builtin.apt: { name: "{{ item }}", state: present }
  loop: "{{ packages }}"
\end{lstlisting}
\subsection{Ad-hoc, Modules \& Errors}
\textbf{Ad-hoc} (one module, no playbook): \texttt{ansible -i inv -m <module> -a '<args>' <target>}. The \texttt{setup} module gathers \textit{facts} (what "Gathering Facts" runs before tasks). \texttt{command} vs \texttt{shell}: \texttt{shell} allows pipes/redirects/vars; \textit{neither is idempotent}. \texttt{ansible-inventory -i inv --list} verifies the inventory; \texttt{--check} is a dry run. A playbook stops at the first failed task unless \texttt{ignore\_errors: true}.
\begin{lstlisting}[language=yaml]
- name: validate
  ansible.builtin.command: httpd -t
  register: out
  changed_when: out.failed
  failed_when: false
- ansible.builtin.file: { path: /data, state: directory } # file|link|touch|absent
- ansible.builtin.lineinfile: { path: /etc/x.conf, line: "Include sites/*" }
\end{lstlisting}
\subsection{Vault, Collections, Roles \& Tags}
\textbf{Vault}: \texttt{ansible-vault create|edit|view|encrypt foo.yml}; \texttt{encrypt\_string 'secret' --name 'pw'} inlines a single var. Pass via \texttt{--vault-id @prompt|/path/file|name@script}. Ciphertext safe to commit. \\
\textbf{Collections} are bundles of modules/roles/plugins. FQCN \texttt{ns.coll.mod} (e.g. \texttt{cisco.ios.ios\_banner}); independent versioning. Install: \texttt{ansible-galaxy collection install <n>} or via \texttt{requirements.yml} (\texttt{collections:} + \texttt{roles:}) then \texttt{install -r}. Sources: \textbf{Galaxy} (community), \textbf{Automation Hub} (RedHat, certified). \\
\textbf{Roles}: \texttt{ansible-galaxy init <n>} creates dirs \texttt{tasks/ handlers/ defaults/ vars/ files/ templates/ meta/ tests/}; use via \texttt{roles:} list. \\
\textbf{Tags} can be used to execute a subset of tasks instead of the whole playbook. Run only specific tags by appending \texttt{--tags <name>} at the end of the ansible-playbook command. There are also two special commands: Tag \textit{always} runs every time, except when explicitly skipped: \texttt{--skip-tags=always}. Tag \textit{never} does not run unless specified with \texttt{--tags=never}
\subsection{Jinja2 Templating}
Rendered on \textbf{controller}, only result sent to target. \texttt{\{\{ inventory\_hostname \}\}} = current target. Filters: \texttt{ipaddr}/\texttt{ipv4}/\texttt{ipv6} (validate/version), \texttt{hash('sha1')}, \texttt{default(x)}, \texttt{min}/\texttt{max}/\texttt{shuffle}.
\begin{lstlisting}[language=yaml]
{% for host in groups['webservers'] %}
BalancerMember http://{{ host }}:{{ port | default(80) }}
{% endfor %}
\end{lstlisting}
